The following finite-group theorem provides the distinction between colors
that the recurrence argument will preserve. Its double-coset assertion is
\cite[Lemma~5.2 and Section~6]{MPNF}. We include the proof, together with
the additional assertion about abelian subgroups, in
Appendix~\ref{sec:finite-groups}.

\begin{theorem}[Finite simple-group obstruction]\label{fin:obstruction}
Let $F$ be a finite nonabelian simple group. There is a nonempty proper
conjugacy-invariant subset $\mathcal S\subset F$, with $1\notin\mathcal S$,
having the following properties.
\begin{enumerate}[label=\textup{(\roman*)}]
\item If $W\in\mathcal S$ and $Z\notin\mathcal S$ commute, then there
exists $B\in F$ such that
\[
 ZW\notin C_F(B)(ZW^{-1})C_F(BZ^2).
\]
\item For every abelian subgroup $A\subset F$ with
$A\cap\mathcal S\ne\varnothing$, the complement $A\setminus\mathcal S$
is a subgroup of $A$.
\end{enumerate}
In particular, $\mathcal S$ is invariant under inversion.
\end{theorem}
